\documentclass[11pt]{amsart}
\usepackage{amsmath,amssymb,amsthm}
\usepackage[margin=1.2in]{geometry}
\usepackage{booktabs}
\usepackage{hyperref}

\newtheorem{theorem}{Theorem}[section]
\newtheorem{lemma}[theorem]{Lemma}
\newtheorem{proposition}[theorem]{Proposition}
\newtheorem{corollary}[theorem]{Corollary}
\newtheorem{conjecture}[theorem]{Conjecture}
\theoremstyle{definition}
\newtheorem{definition}[theorem]{Definition}
\newtheorem{remark}[theorem]{Remark}

\newcommand{\PF}{\mathcal{P}}
\newcommand{\ZZ}{\mathbb{Z}}

% Updated version (August 2026) of the 2016--2019 draft
% "A computational approach to the enumeration of G-parking functions".
% TODO(Tyler): confirm author list and affiliations before circulating.

\title[Enumeration of maximal $G$-parking functions]{A computational
approach to the enumeration of maximal $G$-parking functions, revisited}
\author{Tyler Aden}
\author{Brian Benson}
\date{Updated draft of August 19, 2026}

\begin{document}
\maketitle

\begin{abstract}
Given a connected multigraph $G$ with root $q$, we give an algorithm that
lists every maximal $G$-parking function exactly once, with no duplicate
suppression, by enumerating the lexicographically minimal Dhar burn orders.
Its correctness rests on a structure theorem of independent interest: the
fibers of the map assigning to each $G$-parking function its lex-minimal
burn order are coordinate boxes which partition the set of all $G$-parking
functions, with box tops exactly the maximal parking functions. One search
pass therefore also computes the number of spanning trees $T_G(1,1)$ and the
full degree enumerator $T_G(1,y)$. Our implementation enumerates the
$3\,040\,575$ maximal parking functions of the hypercube $Q_4$ in under two
seconds on a workstation, confirming by direct enumeration the value
previously known only through the chromatic polynomial, and computes
$T_{Q_4}(1,y)$, which appears to be new. We record the resulting data,
including the $133$-element catalog for $Q_3$ that motivated this project.
\end{abstract}

\section{Introduction}

Let $G=(V,E)$ be a connected multigraph with a root $q \in V$. A
\emph{$G$-parking function} \cite{ps} (equivalently a \emph{superstable
configuration} of the abelian sandpile model) is a function
$f\colon V\setminus\{q\}\to\ZZ_{\ge0}$ such that every nonempty
$A\subseteq V\setminus\{q\}$ contains a vertex $v$ with
$f(v) < \sum_{u\notin A} a(v,u)$, where $a(v,u)$ is the edge multiplicity.
The number of $G$-parking functions equals the number of spanning trees of
$G$, and many bijections realizing this are known
\cite{cl,cp,bct}. The \emph{maximal} $G$-parking functions (under the
coordinatewise order) all have the same total degree $g = |E|-|V|+1$, are in
bijection with the acyclic orientations of $G$ with unique source at $q$,
and number $T_G(1,0)$, the Tutte evaluation \cite{bct}.

The original (2016) version of this project sought an algorithm to list the
maximal $Q_n$-parking functions in order to study, experimentally, the
formula for the number of spanning trees of the hypercube,
\begin{equation}\label{eq:stanley}
  c(Q_n) \;=\; 2^{2^n-n-1}\prod_{k=1}^{n} k^{\binom nk},
\end{equation}
for which Stanley's \emph{Enumerative Combinatorics} volume~2 states, at
the end of Example~5.6.10, that ``a direct combinatorial proof of this
formula is not known'' \cite{ec2}. Two combinatorial (though not bijective)
proofs were subsequently given by Bernardi \cite{bernardi}; the second
edition of \cite{ec2} (2024, Notes to Chapter~5) records that a
\emph{bijective} proof remains unknown. The present update replaces the
original breadth-first search over all burn orders --- whose running time
grew roughly like $n^n$ and which required cluster hardware and a trie to
deduplicate output already for $Q_4$ --- by a canonical-order search that
produces each maximal parking function exactly once.

Our contributions are:
\begin{enumerate}
\item an enumeration algorithm for maximal $G$-parking functions with no
duplicate suppression and $O(|V|)$ working memory per search path
(Section~\ref{sec:algo}); we are not aware of any published enumeration
algorithm stated for maximal parking functions (poly-delay enumeration of
acyclic orientations with a prescribed unique source \cite{conte} yields
one implicitly through \cite[Thm.~3.1]{bct});
\item the correctness proof via a fiber theorem: the search leaves are
boxes partitioning all of $\PF(G,q)$ (Theorem~\ref{thm:fiber}), which
simultaneously yields $T_G(1,1)$ and $T_G(1,y)$ in the same pass
(Corollary~\ref{cor:tutte}) and constitutes an explicit form of Merino's
$M$-shellability theorem for cographic matroids \cite{merino01} --- see the
companion note in this repository for that discussion;
\item computational results (Section~\ref{sec:data}): direct-enumeration
confirmation of $T_{Q_4}(1,0)=3\,040\,575$ and
$c(Q_4)=42\,467\,328$; the degree enumerator $T_{Q_4}(1,y)$; the catalog
for $Q_3$; and sandpile-group data through $Q_5$.
\end{enumerate}

\section{Preliminaries}

Dhar's burning criterion \cite{dhar}: $f$ is a $G$-parking function iff,
starting from $q$ burnt and repeatedly burning any vertex $v$ whose burnt
neighbors have total edge weight exceeding $f(v)$, every vertex burns. The
process is abelian, so the following is well defined.

\begin{definition}
The \emph{canonical burn order} $\sigma(f)$ of $f\in\PF(G,q)$ is the burn
order obtained by always burning the smallest-labeled burnable vertex, for a
fixed vertex labeling $q=0,1,\dots,n-1$.
\end{definition}

For an ordering $\sigma$ starting at $q$, write $w_k(v)$ for the weight of
edges from $v$ to $\{\sigma_0,\dots,\sigma_{k-1}\}$, and for $v = \sigma_k$
set $\mathrm{top}_\sigma(v) = w_k(v)-1$ and
$\mathrm{bot}_\sigma(v) = \max\{ w_j(v) : j<k,\ \sigma_j > v \}$ (zero if no
such $j$): the \emph{skip watermark}, i.e.\ the burnt-neighbor weight of $v$
the last time a larger vertex was burned in preference to it.

\section{The fiber theorem and the algorithm}\label{sec:algo}

\begin{theorem}\label{thm:fiber}
For $f\in\PF(G,q)$ with $\sigma=\sigma(f)$,
\[
\{g\in\PF(G,q):\sigma(g)=\sigma\}
 = \{g : \mathrm{bot}_\sigma \le g \le \mathrm{top}_\sigma\},
\]
and every vector in the box is a parking function. Consequently the boxes
partition $\PF(G,q)$; the box tops are exactly the maximal parking
functions, each arising once.
\end{theorem}

\begin{proof}
If $\sigma(g)=\sigma$ then at step $k$ burnability gives
$g(\sigma_k)\le\mathrm{top}_\sigma(\sigma_k)$, and lex-minimality forces
every smaller unburnt $u$ to be unburnable at that step, i.e.\
$g(u)\ge w_k(u)$; maximizing over the skip steps of $u$ gives
$g\ge\mathrm{bot}_\sigma$. Conversely, if
$\mathrm{bot}_\sigma\le g\le\mathrm{top}_\sigma$, induction on $k$ shows the
canonical process for $g$ reproduces $\sigma$: $\sigma_k$ is burnable since
$g(\sigma_k)<w_k(\sigma_k)$, and any smaller unburnt $u$ satisfies
$g(u)\ge\mathrm{bot}_\sigma(u)\ge w_k(u)$ because $u$ is skipped at step
$k$. Hence all of $V$ burns: $g\in\PF(G,q)$ with $\sigma(g)=\sigma$.
For the second statement: every edge contributes to
$w_k(\sigma_k)$ exactly once over the run, so
$\sum_v \mathrm{top}_\sigma(v) = |E|-|V|+1 = g$, making each top maximal
\cite[Prop.~2.3, Cor.~3.2]{bct}; the top lies in its own box, so
$\sigma\mapsto\mathrm{top}_\sigma$ is injective, and a maximal $m$ satisfies
$m\le \mathrm{top}_{\sigma(m)}$ with equal degree, hence equality.
\end{proof}

\begin{corollary}\label{cor:tutte}
With $b_m := \mathrm{bot}_{\sigma(m)}$ and $[k]_y = 1+y+\dots+y^{k-1}$,
\[
T_G(1,y)=\sum_{m\ \mathrm{maximal}}\ \prod_{v\ne q}
  \bigl[\,m(v)-b_m(v)+1\,\bigr]_y,
\qquad
T_G(1,1)=\sum_m \prod_{v\ne q}\bigl(m(v)-b_m(v)+1\bigr).
\]
\end{corollary}

\begin{proof}
Group Merino's identity $T_G(1,y)=\sum_{f}y^{\,g-\deg f}$ \cite{merino97}
by the boxes.
\end{proof}

\subsection*{The search}
The algorithm is a depth-first search over burn prefixes. The state keeps
the burnt set, each vertex's current burnt-neighbor weight
$\mathrm{wcnt}(v)$, and its watermark $w_2(v)$. From a prefix, vertex $v$
may be burned iff it is unburnt and $\mathrm{wcnt}(v) > w_2(v)$; burning $v$
sets $w_2(u)\leftarrow\mathrm{wcnt}(u)$ for every smaller unburnt $u$
(their skip watermark) and adds $v$'s edge weights to its unburnt
neighbors. A leaf at depth $n$ emits the box
$[\mathrm{bot},\mathrm{top}]$; in particular its top is one maximal parking
function, and by Theorem~\ref{thm:fiber} each maximal parking function is
emitted exactly once. A cheap dead-end test (an unburnt vertex with
$\mathrm{wcnt}=w_2$ and no unburnt neighbor can never burn) prunes hopeless
branches. Undoing a burn restores $O(\deg v)$ words, so the working memory
per path is $O(|V|)$, the leaves need no deduplication structure, and
disjoint subtrees parallelize trivially. This replaces both the $n^n$-time
Python search (2016) and the MPI/trie C implementation (2019) of the
earlier drafts.

Correctness and the partition property are validated in the implementation
by exhaustive comparison with brute-force Dhar enumeration on $Q_1$--$Q_3$,
on the $4$-vertex multigraph example of the 2016 code (reproducing its four
maximal parking functions), and on dozens of seeded random connected
multigraphs; tree counts are cross-checked against Matrix-Tree determinants
and, for hypercubes, against \eqref{eq:stanley}.

\section{Data for hypercubes}\label{sec:data}

\begin{center}
\begin{tabular}{rrrrr}
\toprule
$n$ & $g$ & maximal PFs $=T(1,0)$ & spanning trees & time \\
\midrule
2 & 1  & 3 & 4 & --- \\
3 & 5  & 133 & 384 & ms \\
4 & 17 & 3\,040\,575 & 42\,467\,328 & 1.5 s \\
5 & 49 & 81\,768\,640\,551\,939\,777 & $\approx 2.08\times10^{19}$ &
  \emph{not listable} \\
\bottomrule
\end{tabular}
\end{center}

The $Q_4$ row is, to our knowledge, the first direct-enumeration
confirmation of the count $3\,040\,575$, previously derived from the
chromatic polynomial of $Q_4$; the $Q_5$ count is from the chromatic
polynomial of $Q_5$ (OEIS A334278) via
$T(1,0) = |[x]\,\chi_{Q_n}(x)|$, and a Knuth random-descent estimator over
$2\times10^8$ samples of our search tree gives
$(7.4\pm0.5)\times10^{16}$, consistent with it. The degree enumerator of
$Q_4$-parking functions (equivalently $T_{Q_4}(1,y)$, the $h$-vector of the
cographic matroid of $Q_4$), computed in the same $1.5$-second pass:
\[
1,\ 15,\ 120,\ 680,\ 3044,\ 11388,\ 36808,\ 104984,\ 267894,\ 616906,
\]
\[
1287688,\ 2436504,\ 4156516,\ 6306364,\ 8278008,\ 8903016,\ 7016817,\
3040575
\]
(degrees $0$ through $17$). For $Q_3$ the row is
$1, 7, 28, 76, 139, 133$. The Smith normal forms of the reduced Laplacians
give the sandpile groups
$K(Q_3)=\ZZ_2\times\ZZ_8\times\ZZ_{24}$,
$K(Q_4)=\ZZ_2^2\times\ZZ_8\times\ZZ_{24}^3\times\ZZ_{96}$, and
$K(Q_5)=\ZZ_2^5\times\ZZ_6\times\ZZ_{24}^4\times\ZZ_{48}
 \times\ZZ_{192}^3\times\ZZ_{960}$,
consistent with all three theorems of Bai \cite{bai}. The full $Q_3$
catalog ($133$ maximal parking functions, with their boxes) accompanies the
implementation.

\section{Remarks and questions}

The 2016 draft asked about the intersections of the down-sets
$\mathrm{dom}(f)$ of maximal parking functions, following
\cite[Cor.~3.3]{bct}. Theorem~\ref{thm:fiber} answers a sharpened form of
that question: the down-sets can be replaced by canonical disjoint boxes,
eliminating inclusion--exclusion altogether. Questions we consider open and
interesting:
\begin{enumerate}
\item Is the box partition, in lexicographic order of the burn orders, an
$M$-shelling in the sense of Chari? (Verified computationally on all our
test graphs; Merino \cite{merino01} proved existence inductively.)
\item The bottom vector $b_m$ is a statistic on acyclic orientations with
unique source that appears to be new; what does it correspond to under the
known bijections to safe spanning trees?
\item The bijective proof of \eqref{eq:stanley} through this machinery ---
see the companion paper on the support-product theorem.
\end{enumerate}

\begin{thebibliography}{99}
\bibitem{bai} H.~Bai, \emph{On the critical group of the $n$-cube}, Linear
Algebra Appl.\ 369 (2003) 251--261.
\bibitem{bct} B.~Benson, D.~Chakrabarty, P.~Tetali, \emph{$G$-parking
functions, acyclic orientations and spanning trees}, Discrete Math.\ 310
(2010) 1340--1353.
\bibitem{bernardi} O.~Bernardi, \emph{On the spanning trees of the
hypercube and other products of graphs}, Electron.\ J.\ Combin.\ 19(4)
(2012), \#P51.
\bibitem{cl} R.~Cori, Y.~Le~Borgne, \emph{The sand-pile model and Tutte
polynomials}, Adv.\ Appl.\ Math.\ 30 (2003) 44--52.
\bibitem{conte} A.~Conte, R.~Grossi, A.~Marino, R.~Rizzi, \emph{Efficient
enumeration of graph orientations with sources}, Discrete Appl.\ Math.\ 246
(2018) 22--37.
\bibitem{cp} D.~Chebikin, P.~Pylyavskyy, \emph{A family of bijections
between $G$-parking functions and spanning trees}, J.\ Combin.\ Theory
Ser.~A 110 (2005) 31--41.
\bibitem{dhar} D.~Dhar, \emph{Self-organized critical state of sandpile
automaton models}, Phys.\ Rev.\ Lett.\ 64 (1990) 1613--1616.
\bibitem{ec2} R.~P.~Stanley, \emph{Enumerative Combinatorics, Volume 2},
Cambridge Univ.\ Press, 1999; 2nd ed.\ 2024.
\bibitem{merino97} C.~Merino L\'opez, \emph{Chip firing and the Tutte
polynomial}, Ann.\ Comb.\ 1 (1997) 253--259.
\bibitem{merino01} C.~Merino, \emph{The chip firing game and matroid
complexes}, Discrete Math.\ Theor.\ Comput.\ Sci.\ Proc.\ AA (2001)
245--256.
\bibitem{ps} A.~Postnikov, B.~Shapiro, \emph{Trees, parking functions,
syzygies, and deformations of monomial ideals}, Trans.\ Amer.\ Math.\ Soc.\
356 (2004) 3109--3142.
\end{thebibliography}

\end{document}
